\writeSectionFrame{\presentationTitle}{Das Problem}
\begin{frame}{Vor lauter Bäumen den Wald nicht sehen}
	\begin{center}
 		\includegraphics[scale=0.22]{\BILDERPFAD/problem1.png}
 	\end{center}
\end{frame}

\note{
	 In einem Spiel läuft der Spieler durch einen Wald. Der Wald besteht aus mehreren
	 hundert Bäumen, die alle gerendert werden müssen und die alle Speicher belegen.
	 Allerdings sehen sich die Bäume ziemlich ähnlich. Im Prinzip verwenden wir
	 nur zwei unterschiedliche Baumarten.
}


\begin{frame}{Vor lauter Bäumen den Wald nicht sehen}
	\begin{center}
 		\includegraphics[scale=0.22]{\BILDERPFAD/problem2.png}
 	\end{center}
\end{frame}

\note{
	 Jetzt müsste jeder einzelne Baum modelliert werden. Und jeder einzelne Baum
	 belegt auch Speicherplatz.  
}

\begin{frame}{Vor lauter Bäumen den Wald nicht sehen}
	\begin{center}
 		\includegraphics[scale=0.22]{\BILDERPFAD/problem3.png}
 	\end{center}
\end{frame}

\note{
	 Einfacher wäre es, wenn es ein Modell gäbe, das alle Informationen über 
	 Bäume enthält. Die einzelnen Baumobjekte enthalten nur noch die Daten, die
	 bei jedem individuellen Baum anders sind, also die Koordinaten. Wir müssen dann
	 nur einen einzigen Baum modellieren bzw. in unserem Beispiel zwei Bäume, damit
	 wir etwas Abwechslung ins Spiel bekommen. 
}